\documentclass[conference]{IEEEtran}
\usepackage{cite}
\usepackage{amsmath,amssymb,amsfonts}
\usepackage{graphicx}
\usepackage{textcomp}
\usepackage{xcolor}
\usepackage{booktabs}
\usepackage{multirow}
\usepackage{url}
\usepackage{balance}

\graphicspath{{../results/paper_figures/}}

\newif\ifidentifiedversion
\ifdefined\identifiedbuild
    \identifiedversiontrue
\else
    \identifiedversionfalse
\fi

\def\BibTeX{{\rm B\kern-.05em{\sc i\kern-.025em b}\kern-.08em
    T\kern-.1667em\lower.7ex\hbox{E}\kern-.125emX}}

\begin{document}

\title{Air Traffic Control Speaker Diarization with Sortformer: Offline and Streaming Evaluation, Calibration, and Lightweight Domain Adaptation}

\ifidentifiedversion
\author{
\IEEEauthorblockN{Shital Pandey}
\IEEEauthorblockA{Embry-Riddle Aeronautical University\\
1 Aerospace Blvd, Daytona Beach, FL 32114\\
pandeys2@my.erau.edu}
\and
\IEEEauthorblockN{Jianhua Liu}
\IEEEauthorblockA{Embry-Riddle Aeronautical University\\
1 Aerospace Blvd, Daytona Beach, FL 32114\\
liu620@erau.edu}
}
\else
\author{
\IEEEauthorblockN{Anonymous Author 1}
\IEEEauthorblockA{Anonymous Institution\\
Anonymous Location\\
anonymous1@example.com}
\and
\IEEEauthorblockN{Anonymous Author 2}
\IEEEauthorblockA{Anonymous Institution\\
Anonymous Location\\
anonymous2@example.com}
}
\fi

\maketitle

\begin{abstract}
Air traffic control (ATC) speaker diarization is difficult because Very High Frequency (VHF) radio is band-limited, local turns are short, and a single recording may involve dozens of speakers. This paper evaluates the Sortformer family for ATC diarization using author-revised data from the ATC0 dataset. The study jointly examines offline inference, streaming inference, lightweight domain adaptation with a frozen encoder, operating-point calibration, checkpoint sensitivity, and difficulty-stratified behavior under a common segment-level protocol. On the evaluation set, the vanilla offline model reaches a diarization error rate (DER) of 36.6\% and a confusion error rate (CER) of 14.4\%. A 1000-step fine-tuned offline model reduces these values to 14.9\% DER and 6.8\% CER. The vanilla streaming configuration lowers CER to 7.7\% but remains limited by false alarm (FA), yielding 24.8\% DER; selecting a threshold on disjoint calibration recordings and transferring it unchanged to the evaluation set reduces streaming DER to 17.4\% and CER to 6.6\%. Additional analyses show that checkpoint choice is non-monotonic, low-latency streaming is primarily FA-limited, and fine-tuning is most beneficial in segments with multiple speakers or rapid turn changes. The main conclusion is that strong ATC diarization with Sortformer is achievable, but only when adaptation, latency, and threshold calibration are treated as coupled design choices rather than isolated implementation details.
\end{abstract}

\begin{IEEEkeywords}
speaker diarization, air traffic control, Sortformer, streaming inference, domain adaptation, VHF radio
\end{IEEEkeywords}

%===================================================================
\section{Introduction}
\label{sec:intro}
%===================================================================

Air traffic control (ATC) communications are safety-critical voice exchanges carried over Very High Frequency (VHF) radio~\cite{wang2024enhancing, pan2024atcsd}. Automated processing of these communications can support transcription, controller--pilot attribution, and downstream ATC analytics~\cite{wang2024enhancing, zuluaga2022bertraffic}. Speaker diarization, namely the task of determining \emph{who spoke when}, is therefore an important upstream component for ATC speech systems.

For downstream ATC processing, the relevant question is not only whether speech is recognized correctly, but also whether each local transmission is assigned to the correct speaker channel~\cite{wang2024enhancing, zuluaga2022bertraffic}.

ATC audio differs substantially from the meeting and telephone speech that dominates diarization benchmarks~\cite{park2022review, pan2024atcsd}. ATC speech is carried over a band-limited VHF channel, reducing high-frequency speaker information, and push-to-talk operation typically yields one-speaker-at-a-time transmissions rather than continuous open-microphone conversation~\cite{wang2024enhancing, pan2024atcsd}. In the author-revised ATC0 data~\cite{godfrey1994atcc} used here, the 16 recordings used in this study contain 21--90 speakers per recording, well beyond the speaker counts typically encountered in corpora such as CALLHOME and DIHARD~\cite{callhome, ryant2021dihard}.

Modern diarization systems have strong results on standard conversational benchmarks, including end-to-end models such as Sortformer~\cite{park2024sortformer} and its streaming extension~\cite{medennikov2025streaming}, clustering-based pipelines such as pyannote~\cite{bredin2023pyannote}, and linear-complexity streaming models such as Long-Form Streaming End-to-End Neural Diarization (LS-EEND)~\cite{liang2024lseend}. Their behavior on ATC radio speech is less well characterized. In this domain, acoustic mismatch, extreme speaker cardinality, local speaker-slot ambiguity, and threshold sensitivity can all change the conclusions that would be drawn from a meeting-style evaluation.

This paper studies the Sortformer family as a unified offline and streaming framework for ATC diarization. The goal is not to provide a broad benchmark across every available architecture. Instead, the paper addresses two practical questions: \textbf{Q1: How do vanilla offline and streaming Sortformer behave on ATC audio before and after lightweight adaptation?} \textbf{Q2: Once the model family is fixed, how much of the remaining gap is explained by latency and threshold choice rather than by the architecture itself?}

All systems are evaluated under one segment-level protocol: 10\,s analysis windows, 5\,s shift, and scoring against the matching Rich Transcription Time Marked (RTTM) segment with a 0.25\,s collar and overlap retained, as summarized in Fig.~\ref{fig:protocol_overview}. Here, a 0.25\,s collar means that boundary errors within 250\,ms of a reference boundary are ignored, while overlap retained means that overlapping reference speech remains in the scoring rather than being discarded. Under that protocol, the vanilla offline model reaches 36.6\% diarization error rate (DER) and 14.4\% confusion error rate (CER), the fine-tuned offline model reduces these values to 14.9\% and 6.8\%, and threshold transfer reduces streaming DER from 24.8\% to 17.4\%. These results motivate a closer look at adaptation, calibration, and local conversational difficulty.

We endeavor to answer these questions with the following work:
\textbf{First}, we present a controlled evaluation of offline and streaming
Sortformer on four evaluation recordings from three ATC facilities under the
common windowing and scoring protocol shown in Fig
.~\ref{fig:protocol_overview}. \textbf{Second}, we study lightweight domain adaptation through frozen-encoder fine-tuning and show how checkpoint choice changes the DER/CER tradeoff; Fig.~\ref{fig:model_overview} summarizes the shared backbone and the frozen-encoder adaptation path. \textbf{Third}, we quantify the latency--accuracy tradeoff for streaming inference and separate threshold calibration effects from intrinsic model behavior by selecting a common threshold on disjoint calibration recordings and transferring it unchanged to the evaluation set. \textbf{Fourth}, we complement aggregate DER with supporting analyses that explain where the gains arise, including difficulty-stratified evaluation, speaker-count analysis, role-conditioned CER, per-recording DER decomposition, and cue-level embedding evidence.

\begin{figure}[!t]
    \centering
    \includegraphics[width=\columnwidth]{fig_protocol_overview.pdf}
    \caption{Segment-level evaluation protocol. Top: evaluation recordings are processed as overlapping 10\,s analysis windows with 5\,s shift. Bottom: each 10\,s hypothesis is scored against the matching RTTM segment; a 0.25\,s collar is ignored at reference boundaries, and overlapping reference speech is retained.}
    \label{fig:protocol_overview}
\vspace{-2mm}
\end{figure}

The remainder of the paper is organized as follows. Section~\ref{sec:related} places the work relative to end-to-end, streaming, and ATC-specific diarization literature. Section~\ref{sec:setup} defines the dataset, split construction, model configurations, and scoring procedure. Section~\ref{sec:results} presents the main quantitative comparisons. Section~\ref{sec:analysis} examines where the gains come from and which conditions remain difficult. Section~\ref{sec:discussion} interprets the results and states the main limitations.

%===================================================================
\section{Related Work}
\label{sec:related}
%===================================================================

End-to-end neural diarization (EEND) reformulates diarization as multi-label prediction over short acoustic frames, typically on the order of 10--40\,ms depending on feature hop size and model subsampling~\cite{fujita2019eend, park2022review}. Because the assignment of speakers to output slots is arbitrary, EEND models are typically trained with permutation-invariant training (PIT), where the loss is computed against the best matching permutation of reference and predicted speaker labels. Sortformer~\cite{park2024sortformer} replaces PIT with a sort-based speaker ordering objective and combines a Neural Error-correcting Self-supervised Transformer (NEST) Fast Conformer encoder~\cite{huang2025nest} with a deep Transformer decoder. Its streaming extension introduces the Arrival-Order Speaker Cache (AOSC), which maintains local speaker state to reduce online ambiguity~\cite{medennikov2025streaming}. LS-EEND uses retention-based sequence modeling to obtain linear temporal complexity for streaming inference~\cite{liang2024lseend}.

Clustering-based diarization remains important in practice. The pyannote pipeline combines neural speech segmentation, speaker embeddings, and agglomerative clustering~\cite{bredin2023pyannote, plaquet2023powerset}. Such systems combine strong embeddings with explicit postprocessing, which makes them a useful comparison family for ATC radio conditions. In ATC, their effectiveness depends on whether speaker-discriminative embeddings remain sufficiently separated under VHF radio conditions.

Compared with meeting and telephone diarization, the ATC-specific literature discussed here is still relatively small and is often coupled with text or role information~\cite{pan2024atcsd, zuluaga2022bertraffic}. Pan et al.~\cite{pan2024atcsd} proposed an ATC diarization network with voice activity detection, speaker separation, and text-based clustering, reflecting the difficulty of relying on acoustics alone under VHF conditions. Zuluaga-Gomez et al.~\cite{zuluaga2022bertraffic} showed that speaker-role and speaker-change information can be extracted from ATC transcripts. These works highlight an important point: ATC diarization can benefit from lexical or role context, but such context should not be assumed to be available in all deployments. The present study therefore focuses on acoustic diarization itself and examines how a contemporary offline/streaming diarization model transfers to ATC speech.

%===================================================================
\section{Experimental Setup}
\label{sec:setup}
%===================================================================

\subsection{Data, Annotation, and Partitions}

We use author-revised data from the Linguistic Data Consortium (LDC) ATC0 dataset~\cite{godfrey1994atcc}. During inspection of diarization outputs on the original annotations, we found timestamp inconsistencies large enough to affect algorithm evaluation. We therefore revised the annotations used in this study before diarization preparation. The broader annotation revision is ongoing, and the author-revised data used here will be released after the full-dataset update is completed. The author-revised data comprise 16 single-channel VHF recordings from three ATC facilities and span approximately 22 hours of audio. Speaker counts range from 21 to 90 per recording.

Ground-truth speaker annotations are derived from the Segment Time Marked (STM) files and exported to Rich Transcription Time Marked (RTTM) files. Each STM segment is mapped to an RTTM region with the same start/end time and speaker identifier; no speaker merging or relabeling is applied for diarization scoring. Controllers are identified by position labels containing hyphens (for example, D1-1), while pilots appear as airline or callsign identifiers. Four recordings are reserved for evaluation: \texttt{dca\_d1\_1}, \texttt{dca\_d2\_2}, \texttt{dfw\_a1\_1}, and \texttt{log\_id\_1}. The prefixes \texttt{dca}, \texttt{dfw}, and \texttt{log} are the ATC0 file prefixes for the three ATC facilities represented in this study. The four evaluation recordings contain 32, 31, 43, and 88 speakers respectively, giving a deliberately varied test set across both site and speaker cardinality.
These four evaluation recordings are held out from all fine-tuning, validation, and calibration subsets.

The remaining 12 recordings are used for training, validation, and calibration subsets. The study uses three disjoint data partitions: a fine-tuning partition, a threshold-calibration partition, and a final evaluation partition. Fine-tuning uses 90\,s training windows and contains 1453 training windows plus 456 validation windows. Evaluation uses 10\,s analysis windows with 5\,s shift, yielding 3989 scored windows in total: 1056 from \texttt{dca\_d1\_1}, 717 from \texttt{dca\_d2\_2}, 781 from \texttt{dfw\_a1\_1}, and 1435 from \texttt{log\_id\_1}. The fine-tuning train, fine-tuning validation, and calibration subsets all use these same 12 non-evaluation recordings, but with different window durations and sampling policies. For threshold selection, we create the calibration subset by taking non-overlapping 10\,s windows from the first 30 minutes of each of those 12 recordings. Using non-overlapping windows avoids overweighting highly correlated local segments during threshold selection. This produces 2172 calibration windows (181 per recording) and is used only for threshold selection, never for the final evaluation results.

Fig.~\ref{fig:protocol_overview} summarizes the 10\,s windowing, 5\,s shift, and segment-level scoring protocol. Table~\ref{tab:data_partitions} gives the corresponding partition sizes.

\begin{table}[t]
\centering
\caption{Data partitions used in this study. Size is the number of windows.}
\label{tab:data_partitions}
\renewcommand{\arraystretch}{1.05}
\footnotesize
\begin{tabular}{@{}lccc@{}}
\toprule
\textbf{Subset} & \textbf{Recordings} & \textbf{Windowing} & \textbf{\# Windows} \\
\midrule
Fine-tuning train & 12 & 90\,s windows & 1453 \\
Fine-tuning val. & 12 & 90\,s windows & 456 \\
Calibration & 12 & 10\,s, no overlap & 2172 \\
Evaluation & 4 & 10\,s, 5\,s shift & 3989 \\
\bottomrule
\end{tabular}
\vspace{-2mm}
\end{table}

\subsection{Model Configurations and Fine-Tuning}

We evaluate three Sortformer configurations chosen to isolate transfer, latency, and lightweight adaptation.

We refer to the released checkpoints evaluated without ATC-specific adaptation as the \emph{vanilla} configurations, to distinguish them from the fine-tuned variant. Pyannote and LS-EEND are discussed in Section~\ref{sec:related} as comparison families but are not evaluated here; the experiments are intentionally scoped to matched offline and streaming Sortformer variants within one architectural framework.

\textbf{Vanilla offline Sortformer} denotes the released four-speaker offline checkpoint from NVIDIA NeMo~\cite{park2024sortformer}, evaluated without ATC-specific fine-tuning. Inference is performed on 10\,s windows with 5\,s shift.

\textbf{Vanilla streaming Sortformer} is the corresponding four-speaker streaming checkpoint~\cite{medennikov2025streaming}, evaluated at its release operating point before ATC-specific threshold calibration. In the main comparison we report the model's ``medium'' latency configuration, corresponding to 10.0\,s latency. We also evaluate the ``low'' and ``ultra-low'' presets, corresponding to 1.04\,s and 0.32\,s latency.

\textbf{Fine-tuned offline Sortformer} starts from the vanilla offline checkpoint, freezes the NEST encoder, and updates only the Transformer and output layers for 1000 training steps. This design was chosen to preserve the broad acoustic priors in the pre-trained encoder weights while allowing the speaker-assignment layers to specialize. Training uses Adam with decoupled weight decay (AdamW), learning rate $10^{-5}$, 100-step warmup, batch size 4, and bfloat16 (bf16) precision. Checkpoints are retained every 200 steps for the checkpoint-sensitivity study.

For cue-level embedding analysis, we use a TitaNet-Large embedding model based on TitaNet~\cite{koluguri2022titanet} to compute embeddings from isolated ATC cues, where a cue denotes one contiguous transmission from a single speaker to a listener. Fig.~\ref{fig:model_overview} summarizes the Sortformer components shared across the studied configurations and the frozen-encoder adaptation setup.

\begin{figure}[!t]
    \centering
    \includegraphics[width=\columnwidth]{fig_model_adaptation_overview.pdf}
    \caption{Sortformer inference and frozen-encoder adaptation. Top: offline and streaming inference share the same Sortformer backbone, while the streaming variant additionally uses the Arrival-Order Speaker Cache (AOSC). Bottom: adaptation starts from the vanilla offline checkpoint, freezes the NEST/Fast Conformer encoder, and updates only the Transformer/output layers on 90\,s windows from the 12 shared non-evaluation recordings.}
    \label{fig:model_overview}
\vspace{-2mm}
\end{figure}

\subsection{Scoring Protocol and Metric Definitions}

Throughout this paper, \emph{segment} refers to the 10\,s analysis window used for scoring. Sortformer emits up to four \emph{speaker slots}, meaning four output channels that represent locally active speakers within a segment rather than persistent speaker identities across an entire recording. Because those slots are local to each 10\,s analysis segment, evaluation is also performed segment by segment rather than after cross-segment identity stitching. Each hypothesis segment is scored against the matching RTTM region using a 0.25\,s collar with overlap retained. Because the 5\,s shift creates overlapping evaluation windows, window-level observations are not treated as independent samples; instead, the error components are aggregated by recording. This protocol measures segment-level diarization quality while avoiding an additional speaker-linking stage that is not part of the evaluated system.

Diarization error rate (DER) is decomposed into false alarm (FA), missed speech (MISS), and confusion error rate (CER), so that $\mathrm{DER}=\mathrm{FA}+\mathrm{MISS}+\mathrm{CER}$. FA measures predicted speech where the reference contains no speech, MISS measures reference speech omitted by the system, and CER measures speech assigned to the wrong speaker slot.

For the streaming calibration analysis, we sweep a common onset/offset threshold from 0.40 to 0.90 in 0.05 increments on the disjoint calibration subset. Threshold calibration here means selecting the posterior threshold that converts streaming probabilities into active-speaker decisions. Threshold selection is based on minimum DER, with FA and then CER used as tie-breakers if needed. The selected threshold is then transferred unchanged to the four evaluation recordings.

The model family supports four output speaker slots. Fig.~\ref{fig:speaker_capacity} shows the corresponding per-window speaker-count distribution as a within-split percentage, so each series sums to 100\%. This normalization makes the train/validation/evaluation shapes directly comparable despite the different window durations and total numbers of windows. What matters for the present protocol is that no 10\,s evaluation window contains more than four labeled speakers. However, the 90\,s fine-tuning manifests do contain windows with more than four labeled speakers: 13.6\% of training windows and 18.9\% of validation windows have at least five speakers. This mismatch is important when interpreting adaptation results on long ATC windows.

\begin{figure}[t]
    \centering
    \includegraphics[width=\columnwidth]{fig_speaker_capacity.pdf}
    \caption{Per-window speaker-count distribution for Fine-tuning train/val.\ (90\,s) and Evaluation (10\,s), normalized so each split sums to 100\%.}
    \label{fig:speaker_capacity}
\end{figure}

%===================================================================
\section{Results}
\label{sec:results}
%===================================================================

\subsection{Overall Comparison}

Table~\ref{tab:overall_results} summarizes the overall metric decomposition, while Table~\ref{tab:per_recording_results} reports DER/CER pairs by recording. The vanilla offline model serves as the unadapted baseline and reaches 36.6\% DER and 14.4\% CER. The uncalibrated vanilla streaming configuration reduces CER to 7.7\% but still yields 24.8\% DER because FA remains high. The 1000-step fine-tuned offline model gives the strongest overall DER at 14.9\% and also reduces CER to 6.8\%. Displayed component values in Table~\ref{tab:overall_results} are rounded independently, so the rounded FA+MISS+CER may differ from rounded DER by 0.1\%.

\begin{table}[t]
\centering
\caption{Overall segment-level diarization results on the four evaluation recordings.}
\label{tab:overall_results}
\renewcommand{\arraystretch}{1.1}
\footnotesize
\begin{tabular}{@{}lcccc@{}}
\toprule
\textbf{System} & \textbf{DER} & \textbf{FA} & \textbf{MISS} & \textbf{CER} \\
\midrule
Vanilla offline & 36.6\% & 21.2\% & 1.0\% & 14.4\% \\
\shortstack[l]{Vanilla streaming\\(10.0\,s, uncalibrated)} & 24.8\% & 16.6\% & 0.5\% & 7.7\% \\
\shortstack[l]{Calibrated streaming\\(10.0\,s, @ 0.75)} & 17.4\% & 5.5\% & 5.4\% & 6.6\% \\
Fine-tuned offline (1000 steps) & 14.9\% & 6.3\% & 1.9\% & 6.8\% \\
\bottomrule
\end{tabular}
\vspace{-2mm}
\end{table}

\begin{table*}[t]
\centering
\caption{Per-recording DER/CER on the evaluation set. Each cell reports DER / CER.}
\label{tab:per_recording_results}
\renewcommand{\arraystretch}{1.05}
\footnotesize
\begin{tabular}{@{}lccccc@{}}
\toprule
\textbf{System} & \textbf{Overall} & \textbf{\texttt{dca\_d1\_1}} & \textbf{\texttt{dca\_d2\_2}} & \textbf{\texttt{dfw\_a1\_1}} & \textbf{\texttt{log\_id\_1}} \\
\midrule
Vanilla offline & 36.6\% / 14.4\% & 20.9\% / 16.8\% & 20.2\% / 17.1\% & 47.4\% / 15.9\% & 44.7\% / 10.1\% \\
Vanilla streaming (10.0\,s, uncalibrated) & 24.8\% / 7.7\% & 43.3\% / 13.7\% & 43.1\% / 10.3\% & 12.4\% / 4.2\% & 15.5\% / 6.1\% \\
Calibrated streaming (10.0\,s, @ 0.75) & 17.4\% / 6.6\% & 25.9\% / 12.0\% & 25.5\% / 8.7\% & 12.1\% / 3.6\% & 13.1\% / 5.2\% \\
Fine-tuned offline (1000 steps) & 14.9\% / 6.8\% & 12.6\% / 9.0\% & 10.7\% / 7.6\% & 22.1\% / 6.3\% & 11.7\% / 5.4\% \\
\bottomrule
\end{tabular}
\vspace{-2mm}
\end{table*}

The per-recording view is important because the vanilla streaming gains are
not uniform across sites. The calibrated streaming row in
Table~\ref{tab:per_recording_results} uses a threshold selected only on
disjoint calibration recordings and transferred unchanged to the evaluation
set. The overall column aggregates all evaluation windows and is therefore
weighted by window count rather than averaged across the four recording columns. CER decreases on all four evaluation recordings relative to the vanilla offline model, but on the two \texttt{dca} recordings FA rises sharply to 29.0\% and 32.5\%, pushing DER above 43\% despite the lower CER. On the \texttt{dfw} and \texttt{log} recordings, both DER and CER improve simultaneously. The fine-tuned offline system is the strongest uncalibrated overall system and reduces CER on all four recordings relative to the vanilla offline baseline. Relative to the vanilla offline baseline, fine-tuning yields CER reductions of 46.3\%, 55.9\%, 60.7\%, and 45.9\% on \texttt{dca\_d1\_1}, \texttt{dca\_d2\_2}, \texttt{dfw\_a1\_1}, and \texttt{log\_id\_1} respectively, computed from the full-precision per-recording CER values rather than the one-decimal table entries.

\subsection{Checkpoint Sensitivity}

Fig.~\ref{fig:checkpoint_sweep} shows performance across the 200, 400, 600, 800, and 1000-step checkpoints from the fine-tuning run. The response is non-monotonic. DER improves sharply from step 200 to step 400, degrades at step 600, and then partially recovers. The best DER occurs at step 400 (14.79\%), about 0.10 percentage points lower than the 1000-step checkpoint (14.90\%), while the best CER occurs at step 1000 (6.76\%), with step 800 nearly tied on both metrics. We report the 1000-step checkpoint as the main adapted offline system because it corresponds to the fixed 1000-step fine-tuning budget and achieves the lowest CER; if selecting a single checkpoint strictly by DER, step 400 is marginally preferable. The result is practically important: checkpoint selection changes the DER/CER balance and should be reported explicitly rather than treated as a minor training detail.

\begin{figure}[t]
    \centering
    \includegraphics[width=\columnwidth]{fig_checkpoint_sweep.pdf}
    \caption{Checkpoint-response curve for the 1000-step fine-tuning run. The left panel shows DER and CER; the right panel shows FA and MISS. Performance is non-monotonic: the best DER occurs at step 400, whereas the best CER occurs at step 1000.}
    \label{fig:checkpoint_sweep}
\end{figure}

\subsection{Latency--Accuracy Tradeoff}

Fig.~\ref{fig:streaming_frontier} shows the latency frontier across the three streaming configurations. Reducing latency from 10.0\,s to 1.04\,s and 0.32\,s increases DER from 24.8\% to 55.1\% and 70.3\%. The main driver is FA, which rises from 16.6\% (10.0\,s) to 46.2\% (1.04\,s) to 61.1\% (0.32\,s). CER changes much less over the same range, from 7.7\% (10.0\,s) to 8.4\% (1.04\,s) to 8.6\% (0.32\,s). Thus, the low-latency penalty is dominated by activity overtriggering rather than by a wholesale collapse of speaker assignment.

\begin{figure}[t]
    \centering
    \includegraphics[width=\columnwidth]{fig_streaming_frontier.pdf}
    \caption{Streaming latency frontier on the four evaluation recordings. The left panel shows DER versus latency, and the right panel shows CER and FA versus latency. Lower latency increases DER sharply, and the increase is driven primarily by FA rather than by a large increase in CER.}
    \label{fig:streaming_frontier}
\end{figure}

\subsection{Threshold Calibration Across Recordings}

The calibration sweep on the 2172-window disjoint calibration subset selects threshold 0.75 by minimum DER. Fig.~\ref{fig:threshold_transfer} shows both the calibration sweep and the unchanged transfer of that threshold to the evaluation set. Relative to the uncalibrated vanilla streaming configuration, the transferred threshold reduces evaluation DER from 24.8\% to 17.4\% and FA from 16.6\% to 5.5\%. CER also improves modestly from 7.7\% to 6.6\%, while MISS rises from 0.5\% to 5.4\%. The transferred threshold therefore removes much of the \texttt{dca}-specific FA inflation without any tuning on the evaluation recordings, but it does so by trading FA for MISS. After transfer, streaming reaches CER comparable to the fine-tuned offline model (6.6\% versus 6.8\%), but its overall DER remains higher because MISS is substantially larger.

\begin{figure}[t]
    \centering
    \includegraphics[width=\columnwidth]{fig_streaming_threshold_transfer.pdf}
    \caption{Streaming threshold sweep on disjoint calibration recordings and unchanged transfer to the evaluation set. The left panel shows the calibration sweep, and the right panel compares uncalibrated and transferred evaluation metrics. Threshold transfer reduces FA substantially and is therefore central to the interpretation of the streaming results.}
    \label{fig:threshold_transfer}
\end{figure}

%===================================================================
\section{Error Analysis and Supporting Evidence}
\label{sec:analysis}
%===================================================================

\balance
\subsection{Difficulty-Stratified Performance}

Fig.~\ref{fig:window_difficulty} stratifies the evaluation windows by number of active speakers, number of speaker changes, and speech occupancy. Here, \emph{difficulty-stratified} means that windows are grouped by local conversational complexity within each 10\,s segment rather than by the total number of speakers in an entire recording. The CER values in the first two panels are computed by aggregating confusion seconds and total reference speech time within each bin, so bins with more reference speech contribute proportionally more than bins with less speech. Within each bin, the fine-tuned model roughly halves confusion in the harder multi-speaker bins: CER falls from 14.96\% to 6.95\% on two-speaker windows and from 26.42\% to 12.55\% on three-speaker windows. The same pattern holds for rapid-turn windows: CER falls from 22.21\% to 9.94\% on windows with three speaker changes. The vanilla streaming model improves CER over the vanilla offline baseline in these bins, but its main weakness is elevated FA during active speech. In the 25--75\% speech-occupancy bins, streaming produces roughly 0.95--1.06\,s of FA per 10\,s window, compared with 0.13--0.16\,s for the fine-tuned model. Silent-window hallucination is not the dominant failure mode: vanilla streaming produces less FA on silent windows than the vanilla offline model.

\begin{figure}[t]
    \centering
    \includegraphics[width=\columnwidth]{fig_window_difficulty.pdf}
    \caption{Difficulty-stratified analysis on the evaluation windows. From left to right, the panels show CER aggregated within each speaker-count bin, CER aggregated within each speaker-change bin, and mean FA versus speech occupancy. The figure shows that fine-tuning helps most in multi-speaker and rapid-turn segments, while vanilla streaming is limited mainly by FA during active speech.}
    \label{fig:window_difficulty}
\end{figure}

\subsection{Per-Recording and Auxiliary Analyses}

Fig.~\ref{fig:der_comparison} shows the DER decomposition before and after
1000-step fine-tuning. The improvement is not isolated to one recording: all
four evaluation recordings improve in CER, and the \texttt{dfw} and
\texttt{log} recordings also improve strongly in DER through lower FA. The
largest absolute gains combine lower CER with lower FA, especially on the
\texttt{dfw} and \texttt{log} pair. MISS increases modestly overall,
indicating a
slightly more conservative operating point after adaptation.

\begin{figure}[t]
    \centering
    \includegraphics[width=\columnwidth]{fig1_der_comparison.pdf}
    \caption{DER decomposition before and after 1000-step fine-tuning for the
     four evaluation recordings. The figure helps explain \emph{why} DER
     improves: fine-tuning reduces CER on all four recordings and lowers FA
     strongly on the \texttt{dfw} and \texttt{log} pair.}
    \label{fig:der_comparison}
\end{figure}

Role-conditioned analysis shows that the improvement is not evenly distributed across controllers and pilots. Controller CER falls by 44.8\%, 55.1\%, 66.2\%, and 43.4\% across the four evaluation recordings, while pilot CER falls by 38.6\%, 43.3\%, 49.3\%, and 39.6\%. These percentages are computed from role-conditioned evaluation-set totals in the recording order of Table~\ref{tab:per_recording_results}. Controllers account for 63.0\%, 57.8\%, 62.5\%, and 54.5\% of labeled speech time on the evaluation set. The larger controller gains are consistent with the hypothesis that controller speech is acoustically more homogeneous within a recording, although the present study does not isolate a causal explanation.

To test whether ATC audio destroys speaker identity information altogether, we extracted TitaNet-Large embeddings from cue-level segments in \texttt{dca\_d1\_1} and \texttt{log\_id\_1}. This embedding analysis is auxiliary evidence and is not used by the Sortformer diarization systems. Fig.~\ref{fig:embeddings} shows the \texttt{dca\_d1\_1} distribution. The intra-minus-inter cosine similarity margin is 0.318 on \texttt{dca\_d1\_1} and 0.429 on \texttt{log\_id\_1}, indicating that speaker-discriminative information survives the VHF channel to a meaningful extent.

\begin{figure}[t]
    \centering
    \includegraphics[width=\columnwidth]{fig3_embedding_similarity.pdf}
\caption{Cue-level TitaNet-Large similarity distributions for \texttt{dca\_d1\_1}. This auxiliary analysis is not part of the diarization system; it assesses whether VHF ATC audio retains speaker-discriminative cues.}
\label{fig:embeddings}
\end{figure}

%===================================================================
\section{Discussion}
\label{sec:discussion}
%===================================================================

The results support four practical conclusions. First, the vanilla offline Sortformer transfers to ATC only partially: missed speech remains low, but confusion and false alarms remain substantial. Second, the streaming architecture improves local speaker assignment, as shown by lower CER on all four evaluation recordings, but its uncalibrated threshold substantially overfires on the two \texttt{dca} recordings. Third, frozen-encoder fine-tuning provides an effective and economical adaptation path, yielding the best overall DER while reducing confusion most strongly in multi-speaker and rapid-turn windows. Fourth, checkpoint choice and streaming latency materially affect the conclusions drawn from the same model family and should therefore be treated as first-class experimental variables.

Several limitations should also be kept in mind. First, the study evaluates segment-level diarization quality rather than a fully stitched long-form speaker-linking task; a separate cross-window identity reconciliation stage could change the absolute DER values. Second, the evaluation set contains only four recordings, so cross-site generalization should be interpreted cautiously. Third, the model supports four output slots, whereas a non-trivial fraction of the 90\,s fine-tuning windows contain more active speakers. Fourth, the calibrated streaming result reflects a single common threshold selected on disjoint calibration recordings; site-specific calibration or adaptive thresholding may yield different operating points. Operational use in safety-critical avionics contexts would require broader validation beyond this controlled protocol.

The comparison between the fine-tuned offline and calibrated streaming systems is also instructive. After threshold transfer, the streaming model reaches slightly lower CER than the fine-tuned offline model, indicating that local speaker attribution can be competitive once thresholding is corrected. The offline model nevertheless remains superior in DER because it controls FA and MISS more consistently across sites. In practice, calibrated streaming may be attractive for local-turn attribution, whereas archive-quality diarization and downstream analytics still benefit from offline adaptation.

The speaker-count analysis further clarifies why short-window evaluation can be substantially easier than whole-recording speaker counts suggest. Long ATC recordings may involve dozens of speakers, but the local conversational state within a 10\,s segment is much simpler. In this setting, the dominant difficulties are rapid speaker changes, short utterances, and threshold stability under band-limited radio audio rather than global speaker cardinality alone. This explains why a four-slot model remains competitive on 10\,s ATC evaluation windows while still being imperfect for longer fine-tuning windows.

The study is intentionally focused on the Sortformer family because it offers matched offline and streaming variants within one architectural framework. Broader comparisons with text-assisted diarization, larger-capacity models, and explicit cross-window speaker linking remain important future directions.

%===================================================================
\section{Conclusion}
\label{sec:conclusion}
%===================================================================

This paper evaluated Sortformer for ATC speaker diarization under matched offline and streaming conditions and showed that both domain adaptation and operating-point selection materially affect performance. Under the common four-recording segment-level protocol, lightweight offline adaptation reduces DER from 36.6\% to 14.9\% and CER from 14.4\% to 6.8\%. Streaming improves local speaker attribution, but its operating point is strongly threshold-sensitive: threshold transfer reduces streaming DER from 24.8\% to 17.4\% and yields CER of 6.6\% without any tuning on the evaluation recordings. Checkpoint sensitivity is non-monotonic, and low-latency degradation is driven primarily by false alarms. Overall, the results indicate that strong ATC segment-level diarization with Sortformer is feasible, but reliable deployment requires joint treatment of adaptation, latency, and calibration.

\section*{Acknowledgements}
\ifidentifiedversion
We gratefully acknowledge Eclipse Aerospace Inc. for their financial support.
\else
We gratefully acknowledge [sponsor] for their financial support.
\fi

\bibliographystyle{IEEEtran}
\bibliography{references}

\end{document}
